% Fichier engendré par outils/generer-tex.py à partir de 02-geometrie-plane.
% Ne pas éditer à la main : tout le texte provient des commentaires du fichier Lean.
\documentclass[11pt,a4paper]{article}

% Compile aussi bien avec pdflatex qu'avec xelatex ou tectonic.
\usepackage{iftex}
\ifPDFTeX
  \usepackage[T1]{fontenc}
  \usepackage[utf8]{inputenc}
\else
  \usepackage{fontspec}
\fi
\usepackage[french]{babel}
\usepackage{amsmath,amssymb,amsthm}
\usepackage[margin=2.5cm]{geometry}
\usepackage[hidelinks]{hyperref}

\theoremstyle{definition}
\newtheorem{definition}{Définition}
\theoremstyle{plain}
\newtheorem{theoreme}{Théorème}
\newtheorem{lemme}[theoreme]{Lemme}

% Renvoi à la déclaration Lean, une ligne après la démonstration, aligné à droite.
\newcommand{\source}[2]{\par\smallskip\noindent\hfill{\footnotesize\href{#1}{\texttt{#2}}}\par\smallskip}

\title{Géométrie plane}
\date{}

\begin{document}
\maketitle
% Les identifiants Lean en machine à écrire ne se coupent pas : on tolère des
% espacements plus lâches plutôt que des lignes qui débordent.
\sloppy

\section{GeometriePlane}

\noindent
Collège \textemdash{} section \og{} Géométrie plane \fg{}.
Le plan est un espace vectoriel euclidien : un point est un vecteur, une droite est
l'ensemble des \texttt{a + t • u}. Les angles sont ceux de Mathlib, non orientés, mesurés en
radians : l'angle plat vaut \(\pi\) et l'angle droit \(\pi\)/2.
Énoncés et démonstrations en français : voir GeometriePlane.tex.

\subsection{Définitions}

\noindent
Toutes les définitions du fichier sont posées ici, avant les démonstrations.

\begin{definition}
La droite passant par \texttt{a} et dirigée par \texttt{u}.
\end{definition}
\source{https://github.com/Commutator-IO/learn-lean/blob/main/cours/01-college/02-geometrie-plane/GeometriePlane.lean\#L21}{GeometriePlane.lean\#L21}

\begin{definition}
La droite passant par deux points.
\end{definition}
\source{https://github.com/Commutator-IO/learn-lean/blob/main/cours/01-college/02-geometrie-plane/GeometriePlane.lean\#L24}{GeometriePlane.lean\#L24}

\begin{definition}
Deux directions sont parallèles lorsqu'elles sont colinéaires.
\end{definition}
\source{https://github.com/Commutator-IO/learn-lean/blob/main/cours/01-college/02-geometrie-plane/GeometriePlane.lean\#L27}{GeometriePlane.lean\#L27}

\begin{definition}
Deux directions sont perpendiculaires lorsque leur produit scalaire est nul.
\end{definition}
\source{https://github.com/Commutator-IO/learn-lean/blob/main/cours/01-college/02-geometrie-plane/GeometriePlane.lean\#L30}{GeometriePlane.lean\#L30}

\begin{definition}
Le milieu d'un segment.
\end{definition}
\source{https://github.com/Commutator-IO/learn-lean/blob/main/cours/01-college/02-geometrie-plane/GeometriePlane.lean\#L33}{GeometriePlane.lean\#L33}

\begin{definition}
Un quadrilatère \texttt{abcd} est un parallélogramme lorsque les vecteurs de deux côtés
opposés sont égaux.
\end{definition}
\source{https://github.com/Commutator-IO/learn-lean/blob/main/cours/01-college/02-geometrie-plane/GeometriePlane.lean\#L37}{GeometriePlane.lean\#L37}

\subsection{Par deux points distincts passe une droite et une seule}

\begin{theoreme}
Les deux points appartiennent à la droite qu'ils définissent.
\end{theoreme}
% démonstration à transcrire (skill transcrire-preuve-lean)
\source{https://github.com/Commutator-IO/learn-lean/blob/main/cours/01-college/02-geometrie-plane/GeometriePlane.lean\#L42}{GeometriePlane.lean\#L42}

\begin{theoreme}
Et c'est la seule : toute droite contenant \texttt{a} et \texttt{b} distincts leur est égale.
\end{theoreme}
% démonstration à transcrire (skill transcrire-preuve-lean)
\source{https://github.com/Commutator-IO/learn-lean/blob/main/cours/01-college/02-geometrie-plane/GeometriePlane.lean\#L46}{GeometriePlane.lean\#L46}

\subsection{Perpendicularité et parallélisme}

\begin{theoreme}
Si deux droites sont parallèles, toute perpendiculaire à l'une est perpendiculaire à
l'autre.
\end{theoreme}
% démonstration à transcrire (skill transcrire-preuve-lean)
\source{https://github.com/Commutator-IO/learn-lean/blob/main/cours/01-college/02-geometrie-plane/GeometriePlane.lean\#L59}{GeometriePlane.lean\#L59}

\begin{theoreme}
Deux droites parallèles à une même droite sont parallèles entre elles.
\end{theoreme}
% démonstration à transcrire (skill transcrire-preuve-lean)
\source{https://github.com/Commutator-IO/learn-lean/blob/main/cours/01-college/02-geometrie-plane/GeometriePlane.lean\#L66}{GeometriePlane.lean\#L66}

\begin{theoreme}
Dans le plan, deux droites perpendiculaires à une même droite sont parallèles entre
elles. La dimension est essentielle : dans l'espace, deux perpendiculaires à une même
droite ne le sont pas.
\end{theoreme}
% démonstration à transcrire (skill transcrire-preuve-lean)
\source{https://github.com/Commutator-IO/learn-lean/blob/main/cours/01-college/02-geometrie-plane/GeometriePlane.lean\#L75}{GeometriePlane.lean\#L75}

\subsection{Le plus court chemin d'un point à une droite est le segment perpendiculaire}

\begin{theoreme}
Si le vecteur joignant \texttt{p} à \texttt{h} est perpendiculaire à la direction de la droite,
alors \texttt{h} est le point de la droite le plus proche de \texttt{p}.
\end{theoreme}
% démonstration à transcrire (skill transcrire-preuve-lean)
\source{https://github.com/Commutator-IO/learn-lean/blob/main/cours/01-college/02-geometrie-plane/GeometriePlane.lean\#L105}{GeometriePlane.lean\#L105}

\subsection{Angles opposés par le sommet, adjacents, supplémentaires}

\begin{theoreme}
Deux angles opposés par le sommet sont égaux : les deux couples de vecteurs sont
opposés.
\end{theoreme}
% démonstration à transcrire (skill transcrire-preuve-lean)
\source{https://github.com/Commutator-IO/learn-lean/blob/main/cours/01-college/02-geometrie-plane/GeometriePlane.lean\#L118}{GeometriePlane.lean\#L118}

\begin{theoreme}
Deux angles adjacents portés par des demi-droites opposées sont supplémentaires :
leur somme vaut l'angle plat.
\end{theoreme}
% démonstration à transcrire (skill transcrire-preuve-lean)
\source{https://github.com/Commutator-IO/learn-lean/blob/main/cours/01-college/02-geometrie-plane/GeometriePlane.lean\#L124}{GeometriePlane.lean\#L124}

\subsection{Caractérisation de la médiatrice}

\begin{theoreme}
Un point est équidistant de deux points si et seulement s'il appartient à leur
médiatrice.
\end{theoreme}
% démonstration à transcrire (skill transcrire-preuve-lean)
\source{https://github.com/Commutator-IO/learn-lean/blob/main/cours/01-college/02-geometrie-plane/GeometriePlane.lean\#L133}{GeometriePlane.lean\#L133}

\subsection{Somme des angles d'un triangle}

\begin{theoreme}
La somme des angles d'un triangle vaut l'angle plat.
\end{theoreme}
% démonstration à transcrire (skill transcrire-preuve-lean)
\source{https://github.com/Commutator-IO/learn-lean/blob/main/cours/01-college/02-geometrie-plane/GeometriePlane.lean\#L141}{GeometriePlane.lean\#L141}

\subsection{Triangle isocèle : angles à la base}

\begin{theoreme}
Dans un triangle isocèle, les angles à la base sont égaux.
\end{theoreme}
% démonstration à transcrire (skill transcrire-preuve-lean)
\source{https://github.com/Commutator-IO/learn-lean/blob/main/cours/01-college/02-geometrie-plane/GeometriePlane.lean\#L148}{GeometriePlane.lean\#L148}

\subsection{Inégalité triangulaire}

\begin{theoreme}
Inégalité triangulaire : le chemin direct est le plus court.
\end{theoreme}
% démonstration à transcrire (skill transcrire-preuve-lean)
\source{https://github.com/Commutator-IO/learn-lean/blob/main/cours/01-college/02-geometrie-plane/GeometriePlane.lean\#L155}{GeometriePlane.lean\#L155}

\begin{theoreme}
Le cas d'égalité caractérise l'alignement, \texttt{b} étant entre \texttt{a} et \texttt{c}.
\end{theoreme}
% démonstration à transcrire (skill transcrire-preuve-lean)
\source{https://github.com/Commutator-IO/learn-lean/blob/main/cours/01-college/02-geometrie-plane/GeometriePlane.lean\#L159}{GeometriePlane.lean\#L159}

\subsection{Théorème de Pythagore, sa réciproque et sa contraposée}

\begin{theoreme}
Théorème de Pythagore et sa réciproque : le carré de l'hypoténuse est la somme des
carrés des deux autres côtés si et seulement si l'angle est droit.
\end{theoreme}
% démonstration à transcrire (skill transcrire-preuve-lean)
\source{https://github.com/Commutator-IO/learn-lean/blob/main/cours/01-college/02-geometrie-plane/GeometriePlane.lean\#L166}{GeometriePlane.lean\#L166}

\begin{theoreme}
Contraposée : si l'égalité des carrés est en défaut, le triangle n'est pas rectangle.
\end{theoreme}
% démonstration à transcrire (skill transcrire-preuve-lean)
\source{https://github.com/Commutator-IO/learn-lean/blob/main/cours/01-college/02-geometrie-plane/GeometriePlane.lean\#L171}{GeometriePlane.lean\#L171}

\subsection{Théorème des milieux et sa réciproque}

\begin{theoreme}
La droite des milieux est parallèle au troisième côté et de longueur moitié.
\end{theoreme}
% démonstration à transcrire (skill transcrire-preuve-lean)
\source{https://github.com/Commutator-IO/learn-lean/blob/main/cours/01-college/02-geometrie-plane/GeometriePlane.lean\#L178}{GeometriePlane.lean\#L178}

\begin{theoreme}
Conséquence sur les longueurs : le segment des milieux mesure la moitié du troisième
côté.
\end{theoreme}
% démonstration à transcrire (skill transcrire-preuve-lean)
\source{https://github.com/Commutator-IO/learn-lean/blob/main/cours/01-college/02-geometrie-plane/GeometriePlane.lean\#L185}{GeometriePlane.lean\#L185}

\begin{theoreme}
Réciproque : la parallèle menée par le milieu d'un côté coupe le deuxième côté en son
milieu. Écrite vectoriellement : le point du côté \texttt{[a c]} dont l'écart au milieu de
\texttt{[a b]} est parallèle à \texttt{(b c)} avec le bon rapport est le milieu de \texttt{[a c]}.
\end{theoreme}
% démonstration à transcrire (skill transcrire-preuve-lean)
\source{https://github.com/Commutator-IO/learn-lean/blob/main/cours/01-college/02-geometrie-plane/GeometriePlane.lean\#L197}{GeometriePlane.lean\#L197}

\subsection{Théorème de Thalès et sa réciproque}

\begin{theoreme}
Théorème de Thalès : si \texttt{m} et \texttt{n} divisent \texttt{[a b]} et \texttt{[a c]} dans le même rapport
\texttt{t}, alors \texttt{mn} est parallèle à \texttt{bc} et les rapports de longueurs sont égaux à \texttt{t}.
\end{theoreme}
% démonstration à transcrire (skill transcrire-preuve-lean)
\source{https://github.com/Commutator-IO/learn-lean/blob/main/cours/01-college/02-geometrie-plane/GeometriePlane.lean\#L208}{GeometriePlane.lean\#L208}

\begin{theoreme}
Conséquence sur les longueurs : les trois rapports sont égaux.
\end{theoreme}
% démonstration à transcrire (skill transcrire-preuve-lean)
\source{https://github.com/Commutator-IO/learn-lean/blob/main/cours/01-college/02-geometrie-plane/GeometriePlane.lean\#L213}{GeometriePlane.lean\#L213}

\begin{theoreme}
Réciproque de Thalès : si \texttt{mn} est parallèle à \texttt{bc} avec le rapport \texttt{t}, les points
divisent les deux côtés dans ce même rapport.
\end{theoreme}
% démonstration à transcrire (skill transcrire-preuve-lean)
\source{https://github.com/Commutator-IO/learn-lean/blob/main/cours/01-college/02-geometrie-plane/GeometriePlane.lean\#L219}{GeometriePlane.lean\#L219}

\subsection{Relations trigonométriques}

\begin{theoreme}
La relation fondamentale de la trigonométrie.
\end{theoreme}
% démonstration à transcrire (skill transcrire-preuve-lean)
\source{https://github.com/Commutator-IO/learn-lean/blob/main/cours/01-college/02-geometrie-plane/GeometriePlane.lean\#L229}{GeometriePlane.lean\#L229}

\begin{theoreme}
La tangente est le quotient du sinus par le cosinus.
\end{theoreme}
% démonstration à transcrire (skill transcrire-preuve-lean)
\source{https://github.com/Commutator-IO/learn-lean/blob/main/cours/01-college/02-geometrie-plane/GeometriePlane.lean\#L233}{GeometriePlane.lean\#L233}

\subsection{Cercle circonscrit}

\begin{theoreme}
Les médiatrices d'un triangle sont concourantes : il existe un point équidistant des
trois sommets, centre du cercle circonscrit.
\end{theoreme}
% démonstration à transcrire (skill transcrire-preuve-lean)
\source{https://github.com/Commutator-IO/learn-lean/blob/main/cours/01-college/02-geometrie-plane/GeometriePlane.lean\#L239}{GeometriePlane.lean\#L239}

\subsection{Triangle inscrit dans un demi-cercle}

\begin{theoreme}
Un triangle inscrit dans un cercle dont un côté est un diamètre est rectangle, et
réciproquement : c'est le théorème de Thalès sur le cercle.
\end{theoreme}
% démonstration à transcrire (skill transcrire-preuve-lean)
\source{https://github.com/Commutator-IO/learn-lean/blob/main/cours/01-college/02-geometrie-plane/GeometriePlane.lean\#L252}{GeometriePlane.lean\#L252}

\begin{theoreme}
Dans un triangle rectangle, la médiane issue de l'angle droit vaut la moitié de
l'hypoténuse.
\end{theoreme}
% démonstration à transcrire (skill transcrire-preuve-lean)
\source{https://github.com/Commutator-IO/learn-lean/blob/main/cours/01-college/02-geometrie-plane/GeometriePlane.lean\#L258}{GeometriePlane.lean\#L258}

\subsection{Tangente à un cercle}

\begin{theoreme}
La droite passant par un point du cercle et perpendiculaire au rayon en ce point ne
rencontre le cercle qu'en ce point : c'est la tangente.
\end{theoreme}
% démonstration à transcrire (skill transcrire-preuve-lean)
\source{https://github.com/Commutator-IO/learn-lean/blob/main/cours/01-college/02-geometrie-plane/GeometriePlane.lean\#L286}{GeometriePlane.lean\#L286}

\subsection{Parallélogramme}

\begin{theoreme}
Un quadrilatère est un parallélogramme si et seulement si ses diagonales se coupent en
leur milieu.
\end{theoreme}
% démonstration à transcrire (skill transcrire-preuve-lean)
\source{https://github.com/Commutator-IO/learn-lean/blob/main/cours/01-college/02-geometrie-plane/GeometriePlane.lean\#L309}{GeometriePlane.lean\#L309}

\begin{theoreme}
Dans un parallélogramme, les côtés opposés ont la même longueur.
\end{theoreme}
% démonstration à transcrire (skill transcrire-preuve-lean)
\source{https://github.com/Commutator-IO/learn-lean/blob/main/cours/01-college/02-geometrie-plane/GeometriePlane.lean\#L321}{GeometriePlane.lean\#L321}

\begin{theoreme}
Dans un parallélogramme, les angles opposés sont égaux : les deux couples de vecteurs
qui les portent sont opposés.
\end{theoreme}
% démonstration à transcrire (skill transcrire-preuve-lean)
\source{https://github.com/Commutator-IO/learn-lean/blob/main/cours/01-college/02-geometrie-plane/GeometriePlane.lean\#L331}{GeometriePlane.lean\#L331}

\subsection{Repérage : milieu et distance}

\begin{theoreme}
Les coordonnées du milieu sont les moyennes des coordonnées.
\end{theoreme}
% démonstration à transcrire (skill transcrire-preuve-lean)
\source{https://github.com/Commutator-IO/learn-lean/blob/main/cours/01-college/02-geometrie-plane/GeometriePlane.lean\#L347}{GeometriePlane.lean\#L347}

\begin{theoreme}
La distance entre deux points repérés se calcule par le théorème de Pythagore.
\end{theoreme}
% démonstration à transcrire (skill transcrire-preuve-lean)
\source{https://github.com/Commutator-IO/learn-lean/blob/main/cours/01-college/02-geometrie-plane/GeometriePlane.lean\#L354}{GeometriePlane.lean\#L354}

\vfill
\noindent\rule{0.35\linewidth}{0.4pt}\par
\noindent{\footnotesize Leonardo de Moura et Sebastian Ullrich,
\emph{The Lean~4 Theorem Prover and Programming Language},
\emph{Automated Deduction — CADE~28}, Springer, 2021, p.~625--635,
\href{https://doi.org/10.1007/978-3-030-79876-5_37}{doi:10.1007/978-3-030-79876-5\_37}.\par}

\end{document}
